\documentclass[12pt, letterpaper]{article}

% --- Paquetes básicos ---
\usepackage[spanish, es-tabla]{babel}
\usepackage[utf8]{inputenc}
\usepackage[T1]{fontenc}
\usepackage{newpxtext,newpxmath} % Fuente estilo Palatino
\usepackage[margin=3cm]{geometry}
\usepackage{titlesec}
\usepackage[autostyle]{csquotes}

% --- Para dibujar la V de Gowin ---
\usepackage{tikz}
\usetikzlibrary{positioning,babel,calc}

% --- Configuración de la sección (1. IDEA DE...) ---
\titleformat{\section}
  {\normalfont\large\scshape\centering} 
  {\thesection.}
  {0.5em}
  {}

\begin{document}

% --- Encabezado Manual ---
\begin{center}
    \vspace{0.8cm}
    {\Large \bfseries ABOUT PASSWORDS} \\
    \vspace{0.4cm}
\end{center}

\vspace{1.5cm}

\noindent
\textbf{Nivel descriptivo} \newline

\noindent
Título: About Passwords \newline

\noindent
Autor: András Keszthelyi \newline

\noindent
Año: 2013 \newline

\noindent
Nombre del tema: Reflexión práctica sobre el uso, la creación y la seguridad de las contraseñas. \newline

\noindent
Clasificación: Divulgativa y aplicada; el texto presenta observaciones y recomendaciones generales sobre cómo elegir contraseñas más seguras.\newline

\noindent
Resumen en una oración: El autor expone principios básicos para crear contraseñas que sean más resistentes al ataque y, al mismo tiempo, manejables para el usuario. \newline

\noindent
Argumento central: Keszthelyi sostiene que una buena contraseña debe equilibrar la seguridad con la facilidad de recordación, evitando tanto claves demasiado simples como combinaciones imposibles de memorizar.\newline

\noindent
Problemas con el argumento: El texto es útil como orientación general, pero no desarrolla un modelo estadístico o matemático profundo. Por eso, algunas de sus ideas pueden resultar demasiado generales cuando se quiere analizar contraseñas de manera más rigurosa.
\newpage 

\noindent
\textbf{Nivel analítico} \newline

\noindent
Conexión con mi proyecto: Este texto sirve para contextualizar el problema de las contraseñas desde una perspectiva práctica, especialmente porque permite entender por qué los usuarios suelen preferir contraseñas que puedan recordar fácilmente.\newline

\noindent
Lo que NO dice el autor: El autor no construye un modelo formal de predicción ni estudia la distribución de contraseñas dentro de grandes bases de datos. \newline

\noindent
Contraste con otra fuente: A diferencia de investigaciones como las de Ma et al. o Bonneau, este texto no busca medir la seguridad mediante técnicas estadísticas complejas, sino ofrecer una visión más directa sobre el uso cotidiano de contraseñas.   \newline 

\noindent
Evaluación de aplicabilidad: 3. Aporta contexto y observaciones útiles, pero no proporciona herramientas metodológicas fuertes para el análisis cuantitativo del proyecto. \newline

\noindent
Preguntas que le haría al autor: ¿Cómo adaptarían sus recomendaciones a entornos donde se usan gestores de contraseñas o autenticación multifactor? \newline

\noindent
Resumen argumentativo: El texto intenta explicar de manera práctica por qué las personas eligen contraseñas como las eligen. Su valor principal es que ayuda a conectar la teoría de la seguridad con el comportamiento real de los usuarios, aunque no es suficiente por sí solo para un análisis estadístico profundo.

\end{document}
